\chapter{Project Management and Development Evidence}
\label{ch:mgmt}
\section{Work Completed}
\begin{itemize}
\item Next.js application with manager, employee, and shared pages.
\item Role-aware access with signed session cookies and an internal secret header.
\item Trained-role/ByteTrack foot-point preview and provisional floor-distance CLI: editable calibration, motion/reliability checks, metric smoothing and diagnostic exports (since descoped).
\item Two-stage interaction pipeline: trained role detector, BoT-SORT tracking, body-height proximity with a pass-by filter, and a YOLO pose judge; plus a benchmark for detector, interaction and speed KPIs.
\item Executed 17 historical focused floor-distance checks; 44 current interaction unit tests and 5 handoff tests pass; detector benchmarked on 455 held-out images and the interaction pipeline run on 11 local clips.
\end{itemize}

\section{Work Remaining}
Integrating the vision agents with the dashboard, labelled detector/interaction evaluation, the Insight Narration agent on real data, and hardening. Details are in Chapter~\ref{ch:fyp2plan}.

\section{Individual Responsibilities}
\begin{table}[H]
\centering
\caption{Responsibilities}
\begin{tabular}{P{0.25\textwidth}P{0.62\textwidth}}
\toprule
Member & Responsibility\\
\midrule
Moiz Ansari (23I-0523) & Frontend; distance between people; head, hips and feet band handling when the full-person position is rejected\\
Saqib Atif (23I-0769) & Customer/employee pose detection and testing\\
Ruhab Ahmad (23I-0559) & Frontend, report and classification\\
\bottomrule
\end{tabular}
\end{table}

\section{Contribution of Each Team Member}
The following contributions were reported by the team: Moiz worked on frontend development, person-to-person distance and the head/hips/feet band approach for rejected full-person positions; Saqib worked on customer/employee pose detection and testing; Ruhab worked on frontend development, the report and classification. These team-reported areas are distinct from independently verified feature-completion claims. The head/hips/feet categories describe body regions used for distance estimation, not employee/customer role classes.

Git records attribute the provisional floor-distance implementation (\texttt{4a9023a}) and report addition (\texttt{80fb3fd}) to Moiz Ansari. They attribute the interaction detector (\texttt{e3fe2af}) and video-verification record (\texttt{6c25090}) to Ruhab Ahmad. The Git author \texttt{msaqibatifj} is recorded for temporal role classification (\texttt{882f7dc}), the shared OpenVINO model (\texttt{e2b71a5}) and moving-camera engagement (\texttt{face2b5}). The responsibility table above records the team-supplied distribution. Git author names are retained as recorded; the account alias is not used to infer unconfirmed exclusive ownership. These are commit-author records, not a claim that each author exclusively performed every change or a complete account of non-code work.

\section{Repository and Version Control Evidence}
The project is kept in the Git repository \path{muhammad-moiz-ansari/retail-store-analysis} on GitHub. Work is organised on feature branches, and this report corresponds to the \path{feature/interaction} branch, inspected on 20 October 2026 at commit \texttt{6560403}.

\section{Milestones and Commit Evidence}
The actual floor-distance implementation commit is \texttt{4a9023a} (19 October 2026, Moiz Ansari), titled ``feat(vision): add provisional floor-distance estimation pipeline''. A later commit removed the calibrated floor-distance pipeline (\texttt{8c3a6fe}, 20 October), attributed to Git author \texttt{msaqibatifj}. Therefore historical implementation evidence must be read with its commit context. The on-screen interaction flag and one-record-per-interaction output were committed on \path{feature/interaction-flag} (\texttt{20db4c6}, \texttt{ee7947f}, \texttt{0411408}, 19 October 2026, Git author \texttt{msaqibatifj}); the pose judge and benchmark were developed there and ported to \path{feature/interaction} in commit \texttt{6560403} (20 October 2026), after unused code was removed in \texttt{cffa30b}. No Git tags were present in the inspected repository. Selected commit and author-count evidence is reproduced in Appendix~\ref{app:repo}; no project-board or supervisor milestone evidence is inferred from commit dates.

\section{Tools and Tool-Assisted Development}
Development tools include Node.js and npm, a Python virtual environment, Git and GitHub, and the libraries listed in Chapter~\ref{ch:iter1}. Codex assisted with the floor-distance implementation, diagnostic execution and report-source revision. Verification included focused automated tests, actual local trained-model/video runs, inspection of preview/grid outputs, export readability and record-invariant checks. The team also reports using Claude as an AI/development assistant. Codex and Claude are disclosed as development aids; their use does not replace the separately recorded test and review evidence. No specific Claude-generated contribution is attributed to an individual without further team information. Claude Code assisted with porting the pose judge and benchmark to the main branch, the batch and benchmark runs, and the 20 October report revision; outputs were checked with automated tests, a byte-identical comparison of pipeline outputs before and after the port, and inspection of annotated frames.

\section{External Components and Dependencies}
\begin{table}[H]
\centering
\caption{External components}
\begin{tabular}{P{0.25\textwidth}P{0.62\textwidth}}
\toprule
Component & Use\\
\midrule
YOLO implementation \cite{ultralytics} & YOLOv8n pose model as interaction judge; standard YOLO detection for customer/employee classes\\
OpenVINO & CPU inference runtime for the role detector\\
ByteTrack \cite{zhang2022bytetrack} / BoT-SORT & Per-camera tracking (BoT-SORT with appearance re-identification in the interaction pipeline); \texttt{lap} supplies assignment support\\
OpenCV / NumPy \cite{opencvHomography} & Image processing and video input/output; homography fitting, projection and ORB matching in the descoped floor prototype\\
ORB \cite{rublee2011orb} and RANSAC \cite{fischler1981ransac} & Reference-alignment evidence and robust fitting in the descoped floor prototype\\
FastAPI \cite{fastapi} & Application and camera API\\
Next.js and React \cite{nextjs} & Frontend\\
Recharts, Tailwind CSS & Charts, styling\\
SQLite & Application, job and interaction storage\\
\bottomrule
\end{tabular}
\end{table}
\begin{longtable}{P{0.32\textwidth}P{0.55\textwidth}}
\caption{Dependency licence review from installed metadata and upstream notices, 20 October 2026}\\
\toprule
Component & Licence / scope\\
\midrule
\endhead
Ultralytics implementation & AGPL-3.0 or a separate Enterprise licence, according to its publisher \cite{ultralyticsLicense}. Academic use does not by itself remove applicable obligations.\\
Original ByteTrack repository & MIT \cite{bytetrackLicense}; this does not replace the licence of the Ultralytics implementation used by the project.\\
OpenCV & Apache-2.0 for current releases \cite{opencvLicense}.\\
OpenVINO & Apache-2.0 (installed metadata, version 2026.4.1).\\
FastAPI, pytest & MIT (installed distribution metadata).\\
Uvicorn, HTTPX; lap & BSD-3-Clause; BSD-2-Clause for lap (installed metadata).\\
NumPy and PyTorch & BSD-family and bundled-component notices; retain the complete installed distribution notices rather than assuming a single licence covers every component.\\
Pillow & MIT-CMU (installed metadata).\\
Next.js, React, ReactDOM, Recharts, Tailwind CSS & MIT (installed package metadata).\\
Inter, Space Grotesk, JetBrains Mono & SIL Open Font License 1.1 \cite{interLicense,spaceGroteskLicense,jetbrainsMonoLicense}.\\
SQLite & Public-domain software; the application's data rights remain separate.\\
\bottomrule
\end{longtable}
This is a dependency-notice inventory, not a certification of licence compliance or ownership of trained weights/video. Distribution should preserve applicable notices and satisfy the relevant licence terms. Versions installed during testing can differ from manifest ranges; the test environment contained Next.js 16.3.3, React 19.2.8 and Ultralytics 8.4.138.

\section{Risks and Mitigation}
\begin{table}[H]
\centering
\caption{Risk register}
\begin{tabular}{P{0.3\textwidth}P{0.1\textwidth}P{0.4\textwidth}}
\toprule
Risk & Level & Mitigation\\
\midrule
CPU-only inference too slow & High & Frame sampling, small models, batch processing\\
Staff/customer misclassification & High & Zone rules plus appearance, labelled evaluation, manager review\\
Limited footage for evaluation & Medium & Record or label additional clips early in FYP-2\\
Integration slippage & Medium & Bind one agent at a time, keep prototype data fallback\\
Privacy or misuse concerns & Medium & Non-biometric design, decision-support framing\\
Depth ambiguity in a single image & High & Height-ratio check and pose judge; labelled evaluation; reported as a limitation\\
Identity switches on moving footage & High & BoT-SORT re-identification and relinking; count split interactions during evaluation\\
Pose judge too lenient or too strict & Medium & Benchmark with and without the judge; tune facing angle and fraction on labelled clips\\
\bottomrule
\end{tabular}
\end{table}

\section{Security and Privacy Considerations}
Manager routes require a signed HttpOnly session cookie, and direct backend calls require an internal secret header that only the server-side proxy supplies. Secrets come from environment variables. The current login verifies stored account credentials; the earlier arbitrary-email/six-character demo rule no longer describes the tested implementation. The access-control checks are described in Chapter~\ref{ch:testing}.

Floor footage, calibration frames/configuration and generated diagnostics remain local under Git-ignored data directories. The report includes only selected local demonstration frames; no footage upload, persistent person crops or facial-identity records are required by the prototype.

\section{Current Limitations}
\begin{itemize}
\item Vision-derived dashboard metrics still use prototype or demo data; system status is available as a backend endpoint only.
\item The pipeline is offline/batch, not real time.
\item Single-store prototype; no multi-store deployment is claimed.
\item Current authentication uses stored accounts; role-navigation checks cover representative views rather than every application route.
\item The floor-distance prototype was descoped, so no distances in metres are produced; interaction distance is relative to body height and cannot fully resolve depth.
\item Interaction accuracy is measured on one labelled clip only; identity switches on moving footage can split one interaction into several.
\item Interaction results are available through the CLI; dashboard integration remains pending.
\end{itemize}
